\documentclass[11pt]{article}

\usepackage{kotex}
\usepackage{amsmath, amssymb}
\usepackage{booktabs}
\usepackage{graphicx}
\usepackage{hyperref}
\usepackage{geometry}
\usepackage{indentfirst}
\usepackage{float}

\geometry{margin=1in}

\setlength{\parindent}{1em}
\setlength{\parskip}{0pt}

\title{Wordle Solver를 위한 Modified Entropy Maximizing: 파티션 다양성과 균등성의 Trade-off 분석}

\author{
    Useong Jeong
}
\date{}

\begin{document}

\maketitle


\begin{abstract}
본 논문은 Wordle 풀이 문제를 피드백 패턴에 따른 후보 정답 집합의 파티션 문제로 해석하고,
파티션 다양성과 파티션 균등성의 trade-off를 조절하는 Modified Entropy Maximizing (MEM) 방법을 제안한다.
표준 엔트로피 점수는 하나의 추측어가 만드는 피드백 분포의 정보량을 측정하지만,
그 값이 많은 피드백 블록을 생성하는 능력에서 비롯된 것인지,
또는 블록 크기를 균등하게 나누는 능력에서 비롯된 것인지는 명확히 구분하지 않는다.
이에 본 연구에서는 비어 있지 않은 피드백 블록의 수를 반영하는 다양성 점수 $B(g)$와
블록 크기의 균등성을 반영하는 균등성 점수 $U(g)$를 정의하고,
두 요소의 상대적 비중을 조절하는 MEM score $M_{\lambda}(g)$를 제안한다.
특히 $\lambda = 0.5$일 때 제안 점수는 정규화된 표준 엔트로피와 동일한 형태가 되며,
$\lambda$ 값의 변화에 따라 다양성 중심 전략과 균등성 중심 전략을 비교할 수 있다.

실험에서는 2309개의 Wordle 정답 후보와 12947개의 허용 추측어를 사용하여
$\lambda = 0.0$부터 $\lambda = 1.0$까지 solver의 성능을 비교하였다.
그 결과, 표준 엔트로피 baseline인 $\lambda = 0.5$의 평균 시도 횟수는 3.4638이었으며,
파티션 다양성을 더 강조한 설정들에서 전반적으로 더 낮은 평균 시도 횟수가 나타났다.
특히 $\lambda = 1.0$은 평균 시도 횟수 3.4305로 가장 좋은 average-case 성능을 보였다.
반면 $\lambda = 0.7$은 평균 시도 횟수 3.4353으로 근소하게 높았지만,
모든 정답을 5회 이내에 해결하고 4회 이하 해결 비율이 가장 높아 더 안정적인 풀이 분포를 보였다.
이러한 결과는 Wordle solver에서 엔트로피를 구성하는 다양성과 균등성의 trade-off를 명시적으로 조절하는 것이
평균 성능과 worst-case 안정성을 분석하는 데 유용함을 보여준다.
\end{abstract}


\section{Introduction}
\label{sec:introduction}

Wordle은 다섯 글자로 이루어진 숨겨진 단어를 제한된 횟수 안에 맞히는 단어 추측 게임이다.
사용자는 매 시도마다 하나의 단어를 입력하고, 각 글자에 대해 초록색, 노란색, 회색으로 이루어진 피드백을 받는다.
초록색은 해당 글자가 정답의 같은 위치에 있음을 의미하고, 노란색은 해당 글자가 정답에 포함되어 있지만 위치가 다름을 의미한다.
회색은 해당 글자가 정답에 포함되지 않거나, 이미 다른 위치의 피드백에서 사용된 글자 수를 초과했음을 의미한다.
이러한 규칙은 단순해 보이지만, Wordle은 본질적으로 제한된 피드백을 이용하여 가능한 정답 후보를 줄여 나가는 정보 탐색 문제로 해석할 수 있다.

Wordle 풀이 알고리즘에서 중요한 문제는 각 단계에서 어떤 단어를 추측할 것인지 결정하는 것이다.
현재 가능한 정답 후보 집합을 $C$라고 하고, 하나의 추측어를 $g$라고 하자.
각 후보 정답 $a \in C$는 추측어 $g$에 대해 하나의 피드백 패턴 $p(a,g)$를 만든다.
따라서 추측어 $g$는 현재 후보 집합 $C$를 피드백 패턴별 부분집합으로 나누는 역할을 한다.
이를 다음과 같이 나타낼 수 있다.
\[
C_p(g) = \{ a \in C : p(a,g) = p \}.
\]
즉, 하나의 추측어는 현재 가능한 정답 후보들을 여러 개의 피드백 블록으로 분할한다.
좋은 추측어란 실제 피드백을 받은 뒤 남게 되는 후보 집합의 크기를 가능한 한 작게 만드는 추측어라고 볼 수 있다.

이러한 관점에서 널리 사용되는 방법 중 하나가 엔트로피 기반 추측어 선택이다.
추측어 $g$가 만드는 비어 있지 않은 파티션 블록의 크기를 $s_1, s_2, \ldots, s_k$라고 하자.
여기서 $k$는 비어 있지 않은 피드백 블록의 개수이고, $n = |C|$는 현재 후보 정답의 개수이다.
각 피드백 블록이 발생할 확률은 다음과 같이 정의할 수 있다.
\[
q_i = \frac{s_i}{n}.
\]
이때 추측어 $g$의 엔트로피는 다음과 같다.
\[
H(g) = - \sum_{i=1}^{k} q_i \log_2 q_i.
\]
엔트로피가 높다는 것은 해당 추측어가 다양한 피드백 결과를 만들며, 후보들을 비교적 잘 분산시킨다는 의미로 해석할 수 있다.
따라서 표준 엔트로피 기반 solver는 각 단계에서 $H(g)$가 가장 큰 추측어를 선택한다.

그러나 엔트로피 점수는 파티션의 여러 성질을 하나의 값으로 합친다.
예를 들어 어떤 추측어는 매우 많은 피드백 블록을 만들지만, 그중 하나의 블록에 대부분의 후보가 몰릴 수 있다.
반대로 다른 추측어는 생성하는 블록 수는 적지만, 후보들을 비교적 균등하게 나눌 수 있다.
두 경우 모두 엔트로피 값으로 평가할 수는 있지만, 해당 추측어가 좋은 이유가 ``많은 블록을 만들기 때문인지'', 아니면 ``블록 크기를 균등하게 만들기 때문인지''는 명확히 드러나지 않는다.

본 연구에서는 이 점에 주목하여, Wordle 추측어 선택에서 파티션의 품질을 두 가지 요소로 분리한다.
첫 번째 요소는 추측어가 얼마나 많은 비어 있지 않은 피드백 블록을 만드는지를 나타내는 파티션 다양성이다.
두 번째 요소는 생성된 블록들 사이에 후보 정답이 얼마나 균등하게 분포되어 있는지를 나타내는 파티션 균등성이다.
본 연구는 이 두 요소를 결합한 MEM score function을 제안하며, 조절 변수 $\lambda \in [0,1]$를 통해 두 요소 사이의 비중을 조절한다.

제안하는 점수 함수에서 $\lambda = 0.5$인 경우는 기존의 정규화된 엔트로피 점수와 동일하게 해석될 수 있다.
반면 $\lambda > 0.5$인 경우에는 더 많은 피드백 블록을 만드는 추측어를 상대적으로 선호하며, $\lambda < 0.5$인 경우에는 생성된 블록들의 균등성을 더 중요하게 평가한다.
따라서 제안 방법은 기존 엔트로피 기반 접근을 일반화하면서, Wordle 풀이 과정에서 파티션 다양성과 균등성 사이의 trade-off를 명시적으로 조절할 수 있게 한다.

본 논문의 목적은 MEM score function이 Wordle 풀이 성능에 어떤 영향을 주는지 실험적으로 분석하는 것이다.
이를 위해 Wordle 정답 후보 목록과 허용 추측어 목록을 사용하여 여러 $\lambda$ 값에 대해 solver를 실행하고, 평균 시도 횟수, 최대 시도 횟수, 첫 추측어, 시도 횟수 분포 등을 비교한다.
특히 표준 엔트로피에 해당하는 $\lambda = 0.5$를 기준으로, 파티션 다양성 또는 균등성에 더 큰 비중을 두었을 때 성능이 어떻게 달라지는지를 확인한다.

본 논문의 구성은 다음과 같다.
제~\ref{sec:method}장에서는 Wordle 피드백 파티션과 제안하는 MEM score function을 설명한다.
제~\ref{sec:experiments}장에서는 실험에 사용한 데이터셋, 설정, 평가 지표를 제시한다.
제~\ref{sec:results}장에서는 $\lambda$ 값에 따른 실험 결과를 분석한다.
마지막으로 제~\ref{sec:conclusion}장에서는 연구 내용을 요약하고 향후 연구 방향을 제시한다.



















\section{Method}
\label{sec:method}


본 장에서는 Wordle 풀이 과정을 피드백 기반 파티션 문제로 정식화하고, 본 연구에서 제안하는 MEM score function을 설명한다.
제안 방법의 핵심은 하나의 추측어가 만드는 파티션을 단순히 엔트로피 하나로 평가하는 것이 아니라, 파티션의 다양성과 균등성을 분리하여 평가하는 것이다.

\subsection{Wordle Feedback Partitions}

Wordle에서 하나의 추측어는 현재 가능한 정답 후보 집합을 피드백 패턴에 따라 여러 부분집합으로 나눈다.
현재 가능한 정답 후보 집합을 $C$라고 하고, 허용 가능한 추측어 집합을 $G$라고 하자.
어떤 추측어 $g \in G$에 대해, 각 후보 정답 $a \in C$는 하나의 피드백 패턴 $p(a,g)$를 생성한다.

본 연구에서는 Wordle의 피드백을 다음과 같이 인코딩한다.
\[
0 = \text{gray}, \quad
1 = \text{yellow}, \quad
2 = \text{green}.
\]
Wordle은 다섯 글자 단어를 사용하므로 가능한 피드백 패턴의 수는 최대
\[
3^5 = 243
\]
개이다.

추측어 $g$가 주어졌을 때, 동일한 피드백 패턴을 만드는 후보 정답들은 하나의 블록을 이룬다.
피드백 패턴 $p$에 대응되는 블록은 다음과 같이 정의한다.
\[
C_p(g) = \{ a \in C : p(a,g) = p \}.
\]
비어 있지 않은 $C_p(g)$들의 집합은 현재 후보 집합 $C$에 대한 하나의 파티션을 형성한다.

좋은 추측어는 후보 집합을 정보량이 높은 방식으로 나누어야 한다.
즉, 실제 피드백을 받은 뒤 남는 후보 집합의 크기가 작아지도록 현재 후보들을 여러 개의 작은 블록으로 분할하는 추측어가 유리하다.
따라서 Wordle 풀이 문제는 각 단계에서 현재 후보 집합을 가장 효과적으로 분할하는 추측어를 선택하는 문제로 볼 수 있다.

\subsection{MEM Score}

기존의 엔트로피 기반 Wordle solver는 각 추측어가 만드는 피드백 파티션의 정보량을 계산한다.
추측어 $g$에 의해 생성된 비어 있지 않은 파티션 블록의 크기를
\[
s_1, s_2, \ldots, s_k
\]
라고 하자.
여기서 $k$는 비어 있지 않은 피드백 블록의 개수이고, 현재 후보 수는
\[
n = |C|
\]
이다.
각 블록이 발생할 확률은 다음과 같이 정의된다.
\[
q_i = \frac{s_i}{n}.
\]
그러면 추측어 $g$의 엔트로피는 다음과 같다.
\[
H(g) = - \sum_{i=1}^{k} q_i \log_2 q_i.
\]

엔트로피가 클수록 해당 추측어는 후보 집합을 더 정보량 있게 나눈다고 볼 수 있다.
그러나 엔트로피는 파티션의 두 가지 성질을 하나의 값으로 결합한다.
첫 번째는 얼마나 많은 피드백 블록이 생성되는가이고, 두 번째는 생성된 블록들의 크기가 얼마나 균등한가이다.
본 연구에서는 이 두 요소를 각각 파티션 다양성과 파티션 균등성으로 분리한다.

먼저, 파티션 다양성 점수 $B(g)$를 다음과 같이 정의한다.
\[
B(g) =
\frac{\log_2 k}{\log_2(\min(243,n))}.
\]
분모의 $\min(243,n)$은 현재 상태에서 가능한 비어 있지 않은 블록 수의 최댓값을 의미한다.
피드백 패턴은 최대 243개까지 가능하지만, 후보 정답이 $n$개라면 비어 있지 않은 블록은 $n$개를 넘을 수 없다.
따라서 $B(g)$는 현재 상태에서 가능한 최대 블록 수에 비해 추측어 $g$가 얼마나 다양한 피드백 결과를 만들어내는지 나타낸다.

다음으로, 파티션 균등성 점수 $U(g)$를 다음과 같이 정의한다.
\[
U(g) =
\frac{H(g)}{\log_2 k}.
\]
여기서 $\log_2 k$는 $k$개의 블록이 완전히 균등할 때 가능한 최대 엔트로피이다.
따라서 $U(g)$는 생성된 블록들 사이에 후보 정답들이 얼마나 균등하게 분포되어 있는지를 나타낸다.
$U(g)$가 1에 가까울수록 블록 크기가 균등하며, 0에 가까울수록 특정 블록에 후보가 몰려 있음을 의미한다.

본 연구에서 제안하는 MEM score는 $B(g)$와 $U(g)$를 결합하여 정의된다.
\[
M_{\lambda}(g)
=
B(g)^{2\lambda}
\cdot
U(g)^{2(1-\lambda)}.
\]
여기서 $\lambda \in [0,1]$는 파티션 다양성과 파티션 균등성 사이의 상대적 비중을 조절하는 파라미터이다.

$\lambda > 0.5$인 경우에는 $B(g)$의 영향이 더 커지므로, 더 많은 피드백 블록을 생성하는 추측어를 상대적으로 선호한다.
반대로 $\lambda < 0.5$인 경우에는 $U(g)$의 영향이 더 커지므로, 생성된 블록들의 크기를 균등하게 만드는 추측어를 더 선호한다.

특히 $\lambda = 0.5$인 경우, 제안 점수는 기존 정규화 엔트로피와 같은 형태가 된다.
\[
M_{0.5}(g)
=
B(g) \cdot U(g).
\]
이를 전개하면 다음과 같다.
\[
M_{0.5}(g)
=
\frac{\log_2 k}{\log_2(\min(243,n))}
\cdot
\frac{H(g)}{\log_2 k}.
\]
따라서
\[
M_{0.5}(g)
=
\frac{H(g)}{\log_2(\min(243,n))}.
\]
즉, $\lambda = 0.5$는 현재 후보 집합에서 가능한 최대 엔트로피를 기준으로 정규화한 표준 엔트로피 점수에 해당한다.

\subsection{Solver Procedure}

제안한 점수 함수를 사용하는 solver는 각 단계에서 현재 후보 집합 $C$에 대해 가능한 추측어들을 평가하고, 가장 높은 MEM score를 갖는 단어를 선택한다.
전체 절차는 다음과 같다.

\begin{enumerate}
    \item 현재 가능한 정답 후보 집합 $C$를 설정한다.
    \item 각 허용 추측어 $g \in G$에 대해 피드백 파티션을 계산한다.
    \item 각 파티션으로부터 엔트로피 $H(g)$, 다양성 점수 $B(g)$, 균등성 점수 $U(g)$를 계산한다.
    \item MEM score $M_{\lambda}(g)$가 가장 큰 추측어를 선택한다.
    \item 선택한 추측어에 대한 피드백 패턴을 기준으로 후보 집합을 해당 파티션 블록으로 줄인다.
    \item 정답을 찾을 때까지 위 과정을 반복한다.
\end{enumerate}

본 연구의 기본 실험에서는 매 단계에서 전체 허용 추측어 집합 $G$를 평가 대상으로 사용한다.
이는 현재 후보 정답 집합 $C$에 포함되지 않는 단어도 정보 획득을 위한 추측어로 사용할 수 있음을 의미한다.
예를 들어 어떤 단어가 실제 정답 후보는 아니더라도, 남아 있는 후보들을 효과적으로 분할한다면 다음 추측어로 선택될 수 있다.

구현에서는 계산 효율을 높이기 위해 정답 후보와 허용 추측어의 모든 쌍에 대한 피드백 패턴을 미리 계산한 패턴 테이블을 사용한다.
패턴 테이블은 다음과 같이 정의된다.
\[
T[i,j] = p(a_i, g_j),
\]
여기서 $a_i$는 $i$번째 정답 후보이고, $g_j$는 $j$번째 허용 추측어이다.
이 테이블을 사용하면 각 단계에서 피드백 패턴을 반복적으로 새로 계산하지 않고, 현재 후보 집합에 대한 피드백 분포를 빠르게 조회할 수 있다.

또한 동일한 후보 집합이 여러 풀이 경로에서 반복적으로 등장할 수 있으므로, solver는 이미 계산된 상태에 대한 결과를 재사용할 수 있다.
이러한 방식은 여러 $\lambda$ 값에 대해 solver를 반복 실행하는 실험에서 계산 비용을 줄이는 데 도움이 된다.















\section{Experiments}
\label{sec:experiments}

본 장에서는 제안한 MEM 기반 Wordle solver의 성능을 평가하기 위한 실험 설정을 설명한다.
실험의 목적은 MEM score function의 조절 파라미터 $\lambda$가 첫 추측어 선택과 전체 풀이 성능에 어떤 영향을 주는지 분석하는 것이다.
이를 위해 동일한 정답 후보 집합과 허용 추측어 집합을 사용하고, 여러 $\lambda$ 값에 대해 solver를 반복 실행하였다.

\subsection{Dataset}

실험에서는 다섯 글자로 이루어진 Wordle 단어 목록을 사용하였다.
단어 목록은 실제 정답으로 가능한 단어들의 집합과, 사용자가 추측어로 입력할 수 있는 허용 단어들의 집합으로 구분된다.
본 논문에서는 정답 후보 집합을 $A$, 허용 추측어 집합을 $G$로 나타낸다.

실험에 사용한 정답 후보의 개수는 다음과 같다.
\[
|A| = 2309.
\]
허용 추측어 집합은 정답 후보를 포함하도록 구성하였으며, 전체 허용 추측어의 개수는 다음과 같다.
\[
|G| = 12947.
\]
모든 정답 후보는 허용 추측어 집합에도 포함되며,
\[
A \subseteq G
\]
가 성립한다.

각 실험에서는 정답 후보 집합 $A$에 포함된 모든 단어를 한 번씩 실제 정답으로 설정하였다.
즉, solver는 총 $2309$개의 Wordle 문제를 풀게 되며, 각 정답에 대해 몇 번의 추측이 필요한지를 기록하였다.
이 방식은 특정 정답에 대한 성능만을 측정하는 것이 아니라, 전체 정답 후보 집합에 대한 평균적인 solver 성능을 평가할 수 있게 한다.

\subsection{Evaluation Metrics}

solver의 성능은 각 정답을 찾는 데 필요한 추측 횟수를 기준으로 평가하였다.
정답 $a \in A$에 대해 solver가 정답을 맞히는 데 필요한 추측 횟수를 $t(a)$라고 하자.
가장 중요한 평가 지표는 평균 시도 횟수이다.
\[
\text{Mean Attempts}
=
\frac{1}{|A|}
\sum_{a \in A} t(a).
\]
이 값이 작을수록 solver가 평균적으로 더 적은 시도 안에 정답을 찾는다는 의미이다.

또한 최악의 경우 성능을 확인하기 위해 최대 시도 횟수를 함께 측정하였다.
\[
\text{Max Attempts}
=
\max_{a \in A} t(a).
\]
Wordle은 제한된 횟수 안에 정답을 맞히는 게임이므로, 평균 성능뿐만 아니라 특정 정답에서 지나치게 많은 추측이 필요한지도 중요한 평가 기준이 된다.

본 실험에서는 각 $\lambda$ 값에 대해 다음 항목들을 기록하였다.
\begin{itemize}
    \item 선택된 첫 추측어
    \item 평균 시도 횟수
    \item 최대 시도 횟수
    \item 시도 횟수별 정답 개수 분포
    \item 최대 시도 횟수를 발생시킨 정답 목록
\end{itemize}

첫 추측어는 해당 $\lambda$ 값에서 초기 후보 집합 전체를 대상으로 가장 높은 MEM score를 갖는 단어이다.
따라서 $\lambda$ 값이 변하면 첫 추측어도 달라질 수 있다.
이를 통해 $\lambda$가 단순히 후속 탐색 과정뿐만 아니라 초기 정보 탐색 전략에도 영향을 주는지 확인할 수 있다.

\subsection{Lambda Sweep}

본 실험에서는 MEM score function
\[
M_{\lambda}(g)
=
B(g)^{2\lambda}
\cdot
U(g)^{2(1-\lambda)}
\]
에서 $\lambda$ 값을 변화시키며 solver의 성능을 비교하였다.
$\lambda$ 값은 다음과 같이 설정하였다.
\[
\lambda \in \{0.0, 0.1, 0.2, \ldots, 1.0\}.
\]

각 $\lambda$ 값에 대해 solver는 첫 추측어를 고정하지 않고, 초기 후보 집합 $A$와 전체 허용 추측어 집합 $G$를 기준으로 가장 높은 $M_{\lambda}(g)$ 값을 갖는 단어를 첫 추측어로 선택한다.
그 후 피드백에 따라 후보 집합을 줄이고, 각 단계에서 동일한 점수 함수를 사용하여 다음 추측어를 선택한다.

본 실험의 기본 설정에서는 매 단계에서 전체 허용 추측어 집합 $G$를 평가 대상으로 사용하였다.
즉, 현재 가능한 정답 후보 집합에 포함되지 않는 단어라도, 후보들을 효과적으로 분할할 수 있다면 정보 획득을 위한 추측어로 선택될 수 있다.
이 설정은 계산량은 증가하지만, 엔트로피 기반 solver의 정보 탐색 능력을 더 잘 반영한다.

$\lambda = 0.5$인 경우는 중요한 비교 기준으로 사용된다.
앞서 설명한 것처럼 $\lambda = 0.5$일 때 제안한 MEM score는 정규화된 표준 엔트로피 점수와 동일한 형태가 된다.
따라서 본 실험에서는 $\lambda = 0.5$를 표준 엔트로피 기반 solver의 baseline으로 간주한다.

반면 $\lambda < 0.5$인 경우에는 파티션 균등성 $U(g)$에 더 큰 비중을 두고, $\lambda > 0.5$인 경우에는 파티션 다양성 $B(g)$에 더 큰 비중을 둔다.
따라서 $\lambda$ sweep을 통해 표준 엔트로피 기준보다 균등성 또는 다양성을 더 강조했을 때 solver의 평균 성능과 worst-case 성능이 어떻게 달라지는지 분석할 수 있다.
















\section{Results and Discussion}
\label{sec:results}

본 장에서는 $\lambda$ 값에 따른 solver의 성능 변화를 분석한다.
실험은 $\lambda = 0.0$부터 $\lambda = 1.0$까지 0.1 간격으로 수행하였으며,
각 설정에서 선택된 첫 추측어, 평균 시도 횟수, 최대 시도 횟수,
그리고 최대 시도 횟수를 발생시킨 정답의 개수를 기록하였다.

\subsection{\texorpdfstring{$\lambda$}{lambda} 값에 따른 전체 성능 변화}

Table~\ref{tab:lambda-sweep}는 $\lambda$ 값에 따른 solver의 전체 성능을 보여준다.

\begin{table}[H]
\centering
\caption{$\lambda$ 값에 따른 Wordle solver 성능 비교}
\label{tab:lambda-sweep}
\begin{tabular}{c l c c c}
\toprule
$\lambda$ & First guess & Mean attempts & Max attempts & Worst count \\
\midrule
0.0 & roate & 5.0645 & 9 & 4 \\
0.1 & roate & 3.5219 & 5 & 47 \\
0.2 & roate & 3.4937 & 5 & 44 \\
0.3 & roate & 3.4799 & 5 & 42 \\
0.4 & roate & 3.4773 & 5 & 44 \\
0.5 & soare & 3.4638 & 6 & 1 \\
0.6 & reast & 3.4361 & 5 & 53 \\
0.7 & reast & 3.4353 & 5 & 54 \\
0.8 & trace & 3.4309 & 6 & 2 \\
0.9 & trace & 3.4318 & 6 & 2 \\
1.0 & trace & 3.4305 & 6 & 1 \\
\bottomrule
\end{tabular}
\end{table}

Table~\ref{tab:lambda-sweep}에서 먼저 확인할 수 있는 점은 $\lambda$ 값이 첫 추측어 선택에 직접적인 영향을 준다는 것이다.
$\lambda = 0.0$부터 $\lambda = 0.4$까지는 첫 추측어로 \texttt{roate}가 선택되었고,
표준 정규화 엔트로피에 해당하는 $\lambda = 0.5$에서는 \texttt{soare}가 선택되었다.
이후 $\lambda = 0.6$과 $\lambda = 0.7$에서는 \texttt{reast}가 선택되었으며,
$\lambda \geq 0.8$에서는 \texttt{trace}가 첫 추측어로 선택되었다.
이는 $\lambda$가 파티션 균등성과 파티션 다양성 사이의 비중을 조절함에 따라,
초기 후보 집합을 가장 효과적으로 분할한다고 판단되는 단어가 달라짐을 보여준다.

평균 시도 횟수를 기준으로 보면, $\lambda = 1.0$일 때 3.4305로 가장 낮은 평균 시도 횟수를 기록하였다.
이는 표준 엔트로피 baseline인 $\lambda = 0.5$의 평균 시도 횟수 3.4638보다 낮은 값이다.
또한 $\lambda = 0.6$, $\lambda = 0.7$, $\lambda = 0.8$, $\lambda = 0.9$에서도 모두 $\lambda = 0.5$보다 낮은 평균 시도 횟수를 보였다.
따라서 본 실험 설정에서는 표준 엔트로피보다 파티션 다양성에 더 큰 비중을 두는 설정이 평균적인 풀이 성능을 개선하는 경향을 보였다.

반면 $\lambda = 0.0$의 경우 평균 시도 횟수가 5.0645로 크게 증가하였고,
최대 시도 횟수도 9회까지 증가하였다.
$\lambda = 0.0$은 MEM score에서 파티션 균등성만을 강조하는 극단적인 설정이다.
이 경우 생성되는 피드백 블록의 개수는 충분히 고려되지 않으므로,
블록들이 균등하더라도 전체 후보 집합을 충분히 세밀하게 나누지 못할 수 있다.
즉, Wordle solver에서는 단순히 균등한 파티션을 만드는 것뿐만 아니라,
가능한 한 많은 유의미한 피드백 블록을 생성하는 것도 중요하다.

최대 시도 횟수를 기준으로 보면, $\lambda = 0.1$부터 $\lambda = 0.4$까지,
그리고 $\lambda = 0.6$과 $\lambda = 0.7$에서는 최대 시도 횟수가 5회였다.
반면 $\lambda = 0.5$, $\lambda = 0.8$, $\lambda = 0.9$, $\lambda = 1.0$에서는 최대 시도 횟수가 6회였다.
따라서 평균 시도 횟수만 고려하면 $\lambda = 1.0$이 가장 우수하지만,
worst-case 성능까지 함께 고려하면 $\lambda = 0.6$ 또는 $\lambda = 0.7$이 더 균형 잡힌 설정으로 볼 수 있다.
특히 $\lambda = 0.7$은 평균 시도 횟수 3.4353과 최대 시도 횟수 5회를 동시에 달성하므로,
평균 성능과 최악의 경우 성능 사이에서 좋은 trade-off를 보인다.

여기서 Worst count는 각 설정에서 최대 시도 횟수에 도달한 정답의 개수를 의미한다.
따라서 서로 다른 $\lambda$ 값을 비교할 때는 Worst count만 단독으로 해석하기보다,
Max attempts와 함께 해석해야 한다.
예를 들어 $\lambda = 0.5$의 Worst count는 1이지만 최대 시도 횟수는 6회인 반면,
$\lambda = 0.7$의 Worst count는 54이지만 이들은 모두 5회 안에 해결된 경우이다.
Wordle의 제한 횟수를 고려하면, 최대 시도 횟수를 5회로 유지했다는 점에서 $\lambda = 0.7$은 더 안정적인 설정으로 볼 수 있다.

\subsection{시도 횟수 분포와 안정성 분석}

다음으로, 표준 엔트로피 baseline인 $\lambda = 0.5$,
균형 잡힌 설정인 $\lambda = 0.7$,
그리고 평균 시도 횟수가 가장 낮은 $\lambda = 1.0$의 시도 횟수 분포를 비교하였다.
그 결과는 Table~\ref{tab:attempt-distribution-compare}와 같다.

\begin{table}[H]
\centering
\caption{$\lambda = 0.5$, $\lambda = 0.7$, $\lambda = 1.0$의 시도 횟수 분포 비교}
\label{tab:attempt-distribution-compare}
\begin{tabular}{c c c c}
\toprule
Attempts & $\lambda = 0.5$ & $\lambda = 0.7$ & $\lambda = 1.0$ \\
\midrule
1 & 0 (0.00\%) & 0 (0.00\%) & 1 (0.04\%) \\
2 & 45 (1.95\%) & 69 (2.99\%) & 74 (3.20\%) \\
3 & 1213 (52.53\%) & 1220 (52.84\%) & 1232 (53.36\%) \\
4 & 987 (42.75\%) & 966 (41.84\%) & 935 (40.49\%) \\
5 & 63 (2.73\%) & 54 (2.34\%) & 66 (2.86\%) \\
6 & 1 (0.04\%) & 0 (0.00\%) & 1 (0.04\%) \\
\bottomrule
\end{tabular}
\end{table}

Table~\ref{tab:attempt-distribution-compare}에서 볼 수 있듯이,
$\lambda = 0.7$은 세 설정 중 4회 이하로 해결한 정답의 수가 가장 많았다.
전체 2309개의 정답 중 $\lambda = 0.7$은 2255개를 4회 이하에 해결하였으며,
이는 전체의 97.66\%에 해당한다.
반면 표준 엔트로피 baseline인 $\lambda = 0.5$는 2245개,
평균 시도 횟수가 가장 낮은 $\lambda = 1.0$은 2242개를 4회 이하에 해결하였다.
따라서 4회 이하 해결 비율을 기준으로 하면 $\lambda = 0.7$이 가장 안정적인 성능을 보였다고 할 수 있다.

이 결과는 평균 시도 횟수만으로는 드러나지 않는 차이를 보여준다.
$\lambda = 1.0$은 1회, 2회, 3회 안에 해결한 정답 수가 많기 때문에 평균 시도 횟수는 가장 낮다.
그러나 5회 이상이 필요한 정답은 67개로, $\lambda = 0.7$의 54개보다 많다.
즉, $\lambda = 1.0$은 더 많은 정답을 매우 빠르게 해결하는 장점이 있지만,
일부 정답에 대해서는 더 긴 풀이 경로를 만들 수 있다.

반면 $\lambda = 0.7$은 2회와 3회 해결 수를 baseline보다 증가시키면서도,
5회 이상이 필요한 경우를 가장 적게 만들었다.
특히 $\lambda = 0.7$에서는 6회가 필요한 정답이 없었고,
5회 이상 필요한 정답 수도 54개로 세 설정 중 가장 적었다.
Wordle이 최대 6회의 제한을 갖는 게임이라는 점을 고려하면,
4회 이하로 해결되는 정답이 많다는 것은 단순히 평균이 낮다는 것과는 다른 의미의 안정성을 제공한다.
즉, $\lambda = 0.7$은 대부분의 정답을 제한 횟수보다 충분히 이른 시점에 해결하는 경향을 보인다.

전체 시도 횟수로 계산하면, $\lambda = 0.5$에서는 2309개의 정답을 해결하는 데 총 7998회의 추측이 필요하였다.
반면 $\lambda = 0.7$에서는 총 7932회의 추측이 필요하였고,
$\lambda = 1.0$에서는 총 7921회의 추측이 필요하였다.
즉, $\lambda = 0.7$은 표준 엔트로피 baseline에 비해 전체적으로 66회의 추측을 줄였으며,
$\lambda = 1.0$은 77회의 추측을 줄였다.
이는 파티션 다양성을 더 강하게 반영하는 설정이 평균적인 풀이 효율을 개선했음을 보여준다.

결과를 종합하면, MEM score에서 $\lambda$는 solver의 성격을 결정하는 중요한 파라미터이다.
$\lambda$가 너무 작으면 파티션 균등성에 치우쳐 충분한 후보 제거 효과를 얻지 못할 수 있다.
반대로 $\lambda$가 커질수록 더 많은 피드백 블록을 생성하는 추측어를 선호하게 되며,
본 실험에서는 이러한 방향이 평균 시도 횟수를 줄이는 데 효과적이었다.
다만 평균 성능이 가장 좋은 설정이 항상 최대 시도 횟수나 풀이 분포의 안정성 측면에서도 가장 좋은 것은 아니었다.
$\lambda = 1.0$은 평균 시도 횟수 3.4305로 가장 좋은 average-case 성능을 보였지만,
최대 시도 횟수는 6회였고 5회 이상이 필요한 정답 수도 $\lambda = 0.7$보다 많았다.
반면 $\lambda = 0.7$은 평균 시도 횟수 3.4353으로 $\lambda = 1.0$보다 근소하게 높았지만,
4회 이하 해결 수가 가장 많았고, 모든 정답을 5회 이내에 해결하였다.
따라서 $\lambda = 1.0$은 평균 시도 횟수를 최소화하는 설정으로,
$\lambda = 0.7$은 4회 이하 해결 비율과 worst-case 안정성을 함께 고려한 균형 잡힌 설정으로 해석할 수 있다.

































\section{Conclusion}
\label{sec:conclusion}
본 논문에서는 Wordle 풀이 문제를 피드백 패턴에 따른 후보 집합의 파티션 문제로 해석하고,
기존 엔트로피 기반 추측어 선택 방법을 파티션 다양성과 파티션 균등성의 관점에서 확장하였다.
이를 위해 하나의 추측어가 생성하는 비어 있지 않은 피드백 블록의 수를 나타내는 다양성 점수 $B(g)$와,
생성된 블록들의 크기가 얼마나 균등한지를 나타내는 균등성 점수 $U(g)$를 정의하였다.
또한 두 요소의 상대적 중요도를 조절하는 파라미터 $\lambda$를 포함한 MEM score
$M_{\lambda}(g)$를 제안하였다.
이 점수 함수는 $\lambda = 0.5$일 때 정규화된 표준 엔트로피와 동일한 형태가 되며,
$\lambda$ 값을 조절함으로써 기존 엔트로피 기반 solver를 보다 일반적인 형태로 해석할 수 있다.

실험 결과, $\lambda$ 값은 첫 추측어 선택과 전체 풀이 성능 모두에 영향을 주는 중요한 파라미터임을 확인하였다.
표준 엔트로피 baseline에 해당하는 $\lambda = 0.5$에서는 평균 시도 횟수가 3.4638이었으나,
파티션 다양성에 더 큰 비중을 둔 $\lambda > 0.5$ 구간에서는 전반적으로 더 낮은 평균 시도 횟수가 나타났다.
특히 $\lambda = 1.0$은 평균 시도 횟수 3.4305로 가장 좋은 average-case 성능을 보였다.
이는 본 실험 설정에서 후보 집합을 가능한 한 많은 피드백 블록으로 나누는 전략이 평균적인 풀이 효율을 높이는 데 효과적임을 보여준다.

다만 평균 시도 횟수만으로 solver의 성능을 완전히 평가하기는 어렵다.
$\lambda = 1.0$은 가장 낮은 평균 시도 횟수를 기록하였지만, 최대 시도 횟수는 6회였고 일부 정답에서는 비교적 긴 풀이 경로가 발생하였다.
반면 $\lambda = 0.7$은 평균 시도 횟수 3.4353으로 $\lambda = 1.0$보다 근소하게 높았지만,
모든 정답을 5회 이내에 해결하였고 4회 이하 해결 비율도 가장 높았다.
따라서 평균 성능을 최우선으로 고려할 경우에는 $\lambda = 1.0$이 유리하며,
worst-case 안정성과 풀이 분포까지 함께 고려할 경우에는 $\lambda = 0.7$이 더 균형 잡힌 선택으로 해석할 수 있다.

이러한 결과는 Wordle solver에서 단순히 엔트로피 값을 최대화하는 것보다,
엔트로피를 구성하는 요소들을 분리하여 분석하는 것이 의미 있음을 보여준다.
특히 파티션의 균등성만을 지나치게 강조한 $\lambda = 0.0$에서는 평균 시도 횟수와 최대 시도 횟수가 크게 악화되었다.
이는 균등한 파티션을 만드는 것만으로는 충분하지 않으며,
충분히 많은 피드백 블록을 생성하여 후보 집합을 세밀하게 분리하는 능력이 중요하다는 점을 시사한다.
결과적으로 제안한 MEM score는 Wordle 추측어 선택 과정에서 다양성과 균등성 사이의 trade-off를 명시적으로 조절할 수 있는 방법을 제공한다.

향후 연구에서는 몇 가지 방향으로 본 연구를 확장할 수 있다.
첫째, 본 논문에서는 고정된 $\lambda$ 값을 전체 풀이 과정에 동일하게 적용하였지만,
후보 집합의 크기나 풀이 단계에 따라 $\lambda$를 동적으로 조절하는 방법을 고려할 수 있다.
예를 들어 초기 단계에서는 파티션 다양성을 더 강조하고,
후보 수가 줄어든 후반 단계에서는 균등성이나 정답 후보 우선 전략을 더 강하게 반영할 수 있다.
둘째, 평균 시도 횟수뿐만 아니라 실패 확률, 4회 이하 해결 비율, 계산 시간 등을 함께 고려하는 다목적 최적화 관점의 평가도 가능하다.
셋째, tie-breaking 규칙이나 후보 집합 내부 단어를 우선하는 전략을 결합하면,
MEM score의 실제 Wordle 플레이 환경에서의 활용성을 더 높일 수 있을 것이다.

종합하면, 본 연구는 Wordle 풀이에서 엔트로피 기반 접근을 파티션 다양성과 균등성이라는 두 요소로 분해하고,
이를 조절 가능한 점수 함수로 일반화했다는 점에서 의의가 있다.
실험적으로는 표준 엔트로피보다 다양성을 더 강조하는 설정이 평균 성능을 개선할 수 있으며,
특히 $\lambda = 0.7$은 평균 성능과 worst-case 안정성 사이에서 좋은 균형을 보였다.
따라서 제안한 MEM score는 Wordle solver의 성능을 분석하고 조정하기 위한 유용한 기준으로 활용될 수 있다.


\end{document}
